\section{第\ref{sec:sensors}章　マシンの入力}

センサの仕組みは、単一細胞の電流記録、蝸牛の振動の記録、網膜の細胞の計数によって直接測定されている。感度の限界とショットノイズの式は物理から導かれる。センサの符号が情報を最大にするよう調整されているというのは、強い裏づけのある仮説である。

{\footnotesize
\setlength{\tabcolsep}{3pt}\renewcommand{\arraystretch}{1.15}
\begin{longtable}{>{\raggedright}p{0.5cm}>{\raggedright}p{4.0cm}>{\raggedright}p{5.6cm}>{\raggedright}p{2.3cm}>{\raggedright\arraybackslash}p{1.7cm}}
\toprule
No. & 主張 & 実験と測定されたもの & 出典 & 状態 \\
\midrule\endfirsthead
\toprule
No. & 主張 & 実験と測定されたもの & 出典 & 状態 \\
\midrule\endhead
1 & センサは変換器である：受容体が電流を生み、目と耳の感覚細胞の出力は\nt{GLUT}バスへのグルタミン酸である；最初の細胞はスパイクを出さない（図~\ref{fig:transducer}） & カメの桿体と錐体は興奮性アミノ酸を放出する：切り取った膜片のグルタミン酸チャネルがそれを捉える。ラット蝸牛の内有毛細胞：神経線維終末の電流はグルタミン酸AMPA受容体を通り、放出頻度は細胞のなめらかな脱分極とともに $150$\,Hz まで増す & \cite{CopenhagenJahr1989,GlowatzkiFuchs2002} & 測定済み \\
2 & 暗所の光センサは光子1個に約1ピコアンペアの電流で応答する & ヒキガエル桿体の外節への吸引電極：弱い閃光への応答は量子化されており、1事象は $\approx1$\,pA（暗電流の $5\%$）、ばらつき $0.2$\,pA、効率 $0.5$。ヒトは光子1個の到来を偶然より高い頻度で報告する & \cite{Baylor1979,Tinsley2016} & 測定済み \\
3 & 音のセンサは1ナノメートル未満の変位を検出する & メスバウアー法、モルモット蝸牛の基底膜の $16$--$19$\,kHz の位置：聴神経の応答しきい値で膜の速度は約 $0.04$\,mm/s、すなわち変位 $\approx0.35$\,nm（換算は本書による）。測ったのは膜であって、センサの毛束ではない & \cite{Sellick1982,Hudspeth1989} & 測定済み \\
4 & 暗所での精度は光子のショットノイズ~\eqref{eq:photonnoise}で制限される；弱い光では蓄積が長くなる & 式はポアソン統計から導かれる。桿体では弱い光のもとでのノイズが独立な量子事象の和と一致する；明るい背景では閃光への応答が小さく短くなる。「十分の数秒」という数値はここでは別に検証されていない & \cite{Baylor1979} & 理論 \\
5 & 入力のレンジは光で10桁、音で12桁だが、スパイクの発火率は3桁未満 & 音について：特徴周波数での基底膜の応答は $40$--$80$\,dB の帯域で傾き $0.2$\,dB/dB まで増し、圧縮は $100$\,dB を超えても見える。桁数そのものは標準的な推定で、章では出典なし & \cite{Ruggero1997} & 測定済み \\
6 & センサの応答は~\eqref{eq:nakarushton}で、$\sigma$ は平均の明るさを追い、曲線を対数軸に沿ってずらす（図~\ref{fig:adaptation}） & 式は最初、魚の網膜のS電位に当てはめられた。カメの錐体：応答は $V=I V_m/(I+\sigma)$ に従い、明るさの $\sim3.5$ 桁をカバーする；背景に $2$--$3$ 分さらすと、曲線は幅を保ったまま水平にずれる。全体の活動による除算との関係は総説 & \cite{NakaRushton1966,NormannPerlman1979,Carandini2012} & 測定済み \\
7 & センサは変化を伝え、その符号は1スパイクあたりの情報が最大になるよう調整されている（命題~\ref{prop:sensors}） & 原理は理論的に提案された。最も強い証拠：ハエの第1層ニューロンの「コントラスト $\to$ 応答」特性は、自然情景のコントラストの累積分布と一致する。こうして情報が最大になる & \cite{Barlow1961,Laughlin1981} & 仮説 \\
8 & オン型とオフ型のセンサは2線式符号である；イベントカメラはそれをシリコンで再現する & 実験の総説：網膜のオンチャネルとオフチャネルは最初のシナプスですでに分かれ、別々に遮断できる。カメラ：フレームなしの $128\times128$ 画素、レンジ $120$\,dB、遅延 $15$\,\textmu s & \cite{Schiller1992,Lichtsteiner2008,Gallego2022} & 測定済み \\
9 & 目には $\sim10^8$ 個の光センサと $\sim10^6$ 個の出力ニューロンがある：$\sim100$ 分の1への圧縮 & ヒト網膜の全載標本での計数：桿体は平均 $92\cdot10^6$ 個、錐体は $4.6\cdot10^6$ 個；出力（神経節）細胞は目によって $0.7\cdot10^6$ から $1.5\cdot10^6$ 個 & \cite{Curcio1990,CurcioAllen1990} & 測定済み \\
10 & マウスでは目の出力チャネルは30種類を超える & マウス：$11\,000$ 個を超える出力細胞の2光子カルシウムイメージングと応答のクラスタリング。機能的な型は $30$ を明らかに超える。ヒトでは数は測定されていない & \cite{Baden2016} & 測定済み \\
11 & モルモットの目は $\sim875\,000$ bit/s を送る；ヒトの目に換算すると $\sim10^7$ bit/s & モルモット、多電極アレイ、自然情景の中の運動：細胞あたり $6$--$13$ bit/s、$\sim10^5$ 個の出力細胞で $\sim875\,000$ bit/s。ヒト（$\sim10^6$ 個の細胞）の $\sim10^7$ bit/s は著者らの換算 & \cite{Koch2006} & 測定済み \\
12 & 耳はほぼ対数的な周波数マップを持つバンドパスフィルタのバンクである（図~\ref{fig:filterbank}） & 基底膜の振動が最大になる場所は周波数によって決まる；「周波数--場所」関数はほぼ指数関数的で、一つの形でヒトと生きた動物の蝸牛を記述する & \cite{Greenwood1990,Robles2001} & 測定済み \\
13 & 共振器は能動的である：一部のセンサが弱い振動を振幅で数千倍（$66$--$76$\,dB）増幅し、音が大きくなると増幅は下がる & チンチラ蝸牛基部のレーザー振動計測：特徴周波数の弱い純音で、アブミ骨に対する増幅は $66$--$76$\,dB；死後は感度が $60$--$81$\,dB（振幅で $10^3$--$10^4$ 倍）下がり、圧縮が消える。力の源は外有毛細胞 & \cite{Ruggero1997,Robles2001} & 測定済み \\
14 & 耳の増幅器は自励発振の境界に置かれている & 計算：ホップ分岐点の近くでは、レンジの圧縮、鋭い同調、結合音が避けられない。これは聴覚の既知の非線形性のすべてである。蝸牛がまさにそう調整されていることは直接には証明されていない & \cite{Eguiluz2000} & 仮説 \\
15 & 低い周波数ではスパイクは音と位相をそろえて出て（モルモットでは同期は $\sim600$\,Hz から弱まり、$\sim3.5$\,kHz まで認められる）、それから方向が求まる；高い周波数では場所符号が残る & モルモットの聴神経：位相同期は約 $600$\,Hz で落ち始め、$3.5$\,kHz を超えると消える；ネコとサルでは境界がほぼ1オクターブ高い。両耳への到達時間差は総説 & \cite{PalmerRussell1986,Grothe2010} & 測定済み \\
16 & その他のセンサ（表~\ref{tab:sensors}）：嗅覚は $\sim400$ 種類の受容体をニューロンあたり1種類持ち、組み合わせ符号を使う；味覚は一部ATPとセロトニンを使う；触覚は速く順応するセンサと遅く順応するセンサ；温と冷は別々 & マウス：一つの受容体が多くの匂いを認識し、一つの匂いが多くの受容体を活性化し、異なる匂いは異なる組を活性化する。ATP受容体がないと味覚神経は応答しない；味蕾はセロトニンを放出する。ヒトの手の皮膚の単一線維記録：順応の速さで四つの型。冷のチャネルは温のチャネルとは別物 & \cite{Mombaerts2004,Malnic1999,Finger2005,Huang2005,JohanssonVallbo1979,McKemy2002} & 測定済み \\
\bottomrule
\end{longtable}}
